% 需要 \usepackage{booktabs,tabularx}
\begin{table*}[t]
\centering
\footnotesize
\renewcommand{\arraystretch}{1.15}
\setlength{\tabcolsep}{3.0pt}
\caption{主要图级异常检测模型的理论复杂度对比}
\label{tab:model_complexity}
\begin{tabularx}{\textwidth}{lXXXX}
\toprule
模型 & 训练时间复杂度 & 推理时间复杂度 & 参数量 & 激活空间 \\
\midrule
DeepTraLog
& $\mathcal{O}(L(MH+NH^2))$
& $\mathcal{O}(L(MH+NH^2))$
& $\mathcal{O}(T_vDH+LH^2)$
& $\mathcal{O}(M+NH)$ \\
HGT
& $\mathcal{O}(L(N+M)H^2)$
& $\mathcal{O}(L(N+M)H^2)$
& $\mathcal{O}(LRH^2)$
& $\mathcal{O}(MH+NH)$ \\
OCHetGCN
& $\mathcal{O}(L(N+M)H^2)$
& $\mathcal{O}(L(N+M)H^2)$
& $\mathcal{O}(LRH^2)$
& $\mathcal{O}(MH+PH+NH)$ \\
HRGCN
& $\mathcal{O}(2L(N+M)H^2)$
& $\mathcal{O}(L(N+M)H^2)$
& $\mathcal{O}(LRH^2)$
& $\mathcal{O}(MH+PH+NH)$ \\
GLocalKD
& $\mathcal{O}(2L(N+M)H^2)$
& $\mathcal{O}(2L(N+M)H^2)$
& $\mathcal{O}(2LRH^2)$
& $\mathcal{O}(MH+PH+NH)$ \\
HRA-GNN
& $\mathcal{O}(2L(N+M)H^2)$
& $\mathcal{O}(L(N+M)H^2)$
& $\mathcal{O}(LRH^2)$
& $\mathcal{O}(MH+PH+NH+LRH)$ \\
SIGNET
& $\mathcal{O}(\mathcal{G}Ld^2(m+n)+\mathcal{G}nd(d_f+d_f^\ast)+\mathcal{G}Bd)$
& $\mathcal{O}(\mathcal{G}Ld^2(m+n)+\mathcal{G}Bd)$
& $\mathcal{O}(Ld^2+d(d_f+d_f^\ast))$
& $\mathcal{O}(B^2)$ \\
CVTGAD
& $\mathcal{O}(L(MH+NH^2)+(N+B)D_c^2+(N^2+B^2)D_c)$
& $\mathcal{O}(L(MH+NH^2)+(N^2+B^2)D_c)$
& $\mathcal{O}(LH^2+D_c^2)$
& $\mathcal{O}(N^2+B^2)$ \\
MUSE
& $\mathcal{O}(N^2)$
& $\mathcal{O}(N^2)$
& $\mathcal{O}(LH^2)$
& $\mathcal{O}(N^2)$ \\
GLADMamba
& $\mathcal{O}(L(N+M)H^2+(N^2+B^2)D_c)$
& $\mathcal{O}(L(N+M)H^2)$
& $\mathcal{O}(LH^2+D_c^2+D_cS)$
& $\mathcal{O}(N^2)$ \\
\bottomrule
\end{tabularx}
\par\vspace{2pt}
\begin{minipage}{\textwidth}
\scriptsize
\textbf{(1) 系数说明：}
$N$、$M$分别为一个mini-batch中的节点和边总数，$B$为图数量，
$H$为隐藏维度，$L$为层数，$R$为关系类型数，$P\le M$为有效
``目标节点--关系''二元组数，$D_c=LH$，$S$为Mamba状态维度。
SIGNET一行保留原论文符号（$\mathcal{G}$为图总数，$n,m$为平均节点和边数，
$d,d_f,d_f^\ast$为原论文维度参数）。
HRA-GNN训练系数2来自自监督增强分支；HRGCN系数2来自双编码分支；
GLocalKD系数2来自教师--学生双网络。
CVTGAD和GLADMamba的$N^2$项分别来自节点交叉注意力和节点对比损失矩阵；
MUSE的$N^2$来自稠密邻接重构。

\medskip
\textbf{(2) 复杂度分层：}
十个模型按推理时间复杂度可划分为三档：
\textbf{线性档}（DeepTraLog、HGT、OCHetGCN、HRGCN、HRA-GNN、GLADMamba推理）
均为稀疏消息传递，推理复杂度$\mathcal{O}(L(N+M)H^2)$，
对图规模$N+M$近似线性；其中HRA-GNN参数量约为HRGCN的一半
（$\mathcal{O}(LRH^2)$ vs $\mathcal{O}(2LRH^2)$），
因为HRA-GNN不需要额外的拼接投影矩阵。
\textbf{二次档}（CVTGAD、GLADMamba训练）推理仍为线性，但训练中分别因
节点交叉注意力和对比损失矩阵引入$\mathcal{O}(N^2)$项。
\textbf{全二次档}（MUSE）训练和推理均为$\mathcal{O}(N^2)$，
来自稠密邻接重构，不适用于大图。
\end{minipage}
\end{table*}
